\section{What the counterexample changes}\label{sec:limits}

\subsection{For lower bounds and for constructions}\label{sec:limits-machines}

The fifth paper proposed \cref{prob:q91} as a machine-free route to lower bounds: an affirmative answer, for every fixed
$\Delta$, would bound by $1/2+o(1)$ the rate of every finite-alphabet storage code on the actual cone relations of
oblivious linear-time computations~\cite[Proposition~9.2]{DirectedRecovery2026}, and combined with cheap depth reduction
this would exclude oblivious linear-time transposers. Theorems A and~A$'$ close this route in its machine-free form:
bounded degree and the paired-distance condition do not force rate $1/2$, for linear codes over a fixed field or for
binary codes. Corollary U closes the undirected form of the route, the fourth paper's
Question~6.1~\cite{CurvedMotion2026}, in the same way: bounded degree without short odd cycles does not force rate $1/2$
either. The conclusion of the fifth paper's Proposition~9.2 for actual cone relations is not refuted, because nothing here shows that
our relations arise from tape computations. A half-rate bound for the cone relations of linear-time machines, or any bound
below the rates of Theorems A and A$^*$, must therefore use structure of tape computations beyond bounded degree and
paired distances. A bound $1-\rho$ with $\rho>0$ independent of the degree is not excluded by our construction (\cref{sec:open}). The fifth paper
records such structure, for instance its transport bound for matchings of cone
relations~\cite[Lemmas~3.4 and~3.5]{DirectedRecovery2026}. Several later papers of the series state results conditional on an
affirmative answer to \cref{prob:q91}; as implications they remain correct, but their hypothesis is now known to be false in the
generality stated, and they need the native form of the hypothesis instead.

On the construction side, a linear-time transposer needs its cone relations to carry codes whose rate beats the genie
cost~\cite[Section~2]{DirectedRecovery2026}; for genies of cost $\eps L$ this means rate close to $1$. The relations of
Theorem A carry codes of rate above one half, which defeats the half-bound route but not this requirement; they come with
an astronomically large degree, a characteristic-$127$ linear structure and a twist by a character of the congruence
subgroup $\Gamma^+(\lambda_0)$, and no tape realization of them is known. Whether a linear-time oblivious machine can have
cone relations with this kind of twisted, locally separated structure is the question that the counterexample leaves for
the construction side.

\subsection{Limits of the method}\label{sec:limits-method}

The rates guaranteed by Theorem A$^*$ approach $(\sqrt3-1)^2=4-2\sqrt3\approx0.536$. Two features of the construction set
this limit. Our lower bound for the Goren--Oort codes, from Riemann--Roch, is the rate $(1-t/(p^2+1))^2$ of a tensor
product of two local codes, up to lower-order terms; we prove no upper bound beyond the trivial one, but a rate close to
$1$ would need a different source. And the separated structures of \cref{thm:D} have degeneracy at least about
$(2-\sqrt3)(p^2+1)$, since the fluid rule gives every lifted vertex expected load $\tau$ per unit degree; whether other
lift selections can do better is \cref{rem:optimal}.

\begin{remark}[Heuristic: the constant of Theorem D is optimal]\label{rem:optimal}
Heuristic, not used in any proof. The fluid rule balances the loads exactly: every lifted vertex has expected load
$\tau$ per unit degree whatever its label (\cref{lem:fluid}). For an order of the lifted vertices of a random-like graph,
a lifted vertex at a given position has a definite fraction of its neighbours later, and summing the resulting
constraints over all positions suggests that no lift selection of such a graph has degeneracy below
$(2-\sqrt3-o(1))\delta$. If so, constructions of the present type, with tensor-shaped rates, cannot exceed $4-2\sqrt3$.
\end{remark}

\begin{remark}[Congruence levels are never untwisted]\label{rem:untwisted}
This remark is not used elsewhere. In the setting of \cref{thm:A}, no $\Gamma^+(\lambda_0^j)$ is contained in the kernel of
the twist $\chi$. Otherwise the classes would descend to $X_j^1$, every recovery set of $N_j$ would lie in the other class,
and the signed class of the proof of \cref{prop:descent}(c) would separate $L_x$ from $R_x$ in the incidence graph; by the
fifth paper's half bound for such relations~\cite[Lemma~4.1]{DirectedRecovery2026}, every storage code would then have rate
at most $1/2$, contradicting \cref{thm:G}. So the arithmetic forces the twist: the subgroup $\ker\chi$ contains none of
the congruence subgroups $\Gamma^+(\lambda_0^j)$.
\end{remark}

\subsection{Characteristic \texorpdfstring{$p$}{p}}\label{sec:limits-char}

Theorem G is a phenomenon of positive characteristic. Along the congruence tower the normalized dimension of the Tanner
codes stays at least $(k-1)^2/(p^2+1)$, while along towers with $p$-group quotients the normalized dimension of the same Tanner
operator converges to a rank over a division ring, as shown by Jaikin-Zapirain and Souza~\cite[Theorem~8.1]{JaikinSouza2026};
the twelfth paper of the series uses this difference to show that L\"uck approximation fails in positive
characteristic~\cite{Luck2026}. The sixth paper of the series bounds linear storage codes along towers obtained from pro-$p$
embeddings in characteristic~$p$, through the same theorem of Jaikin-Zapirain and Souza~\cite{Covers2026}; our tower has
quotients of order a power of $509$ and codes of characteristic $127$, so it lies outside those results.

\begin{remark}[Heuristic: characteristic zero]\label{rem:char0}
Heuristic, not used in any proof. Over subfields of $\C$, L\"uck's approximation theorem~\cite{Luck1994} and its
extensions~\cite[Section~1 and Theorem~10.1]{JaikinSurvey2019} make
normalized dimensions of Tanner codes converge, along every tower of normal subgroups with trivial intersection, to the
von Neumann dimension of the Tanner code on the tree, which is the same for all towers. For relations that are perfectly
separated on the tree, as ours are, this common limit appears to be at most $1/2$, by comparison with the untwisted
covers, which obey the half bound. This suggests that linear codes in characteristic zero cannot replace the
Goren--Oort codes, and that the half bound may hold for complex-linear codes on all bounded-degree relations with the
paired-distance property; the fifth paper proves it for complex-linear codes with bounded coefficients at fixed vector
width~\cite[Theorem~4.5]{DirectedRecovery2026}.
\end{remark}
